\begin{exercise}{5}
Show that every simplex $x\in X_n$ is of the form $\alpha^*(y)$ for a
surjection $\alpha\colon[n]\to[m]$ and a non-degenerate $m$-simplex $y$,
and show that the pair $(\alpha,y)$ is uniquely determined by
$x$.\footnote{The original has $\alpha\colon[m]\to[n]$ and
``a non-degenerate $n$-simplex $y$''. These should read
$\alpha\colon[n]\to[m]$ and ``a non-degenerate $m$-simplex $y$''.
Since $X$ is contravariant, $\alpha\colon[n]\to[m]$ induces
$\alpha^*\colon X_m\to X_n$, so $y$ must lie in $X_m$ for
$\alpha^*(y)=x\in X_n$.}
% Correction source: IIC/chapters/chA-exercises/sA-exercises.tex,
% label ex:eilenberg-zilber, translator's note, lines 45--48.

\begin{proof}
Existence follows by repeatedly expressing a degenerate simplex as
$s_i^*$ of a simplex of smaller dimension (Definition~1.1.7).
The dimension strictly decreases, so the process ends at a
nondegenerate simplex; the composite of the degeneracy maps in
$\Delta$ is a surjection.

For uniqueness, suppose $x=\alpha^*y=\beta^*z$, where
$\alpha\colon[n]\twoheadrightarrow[p]$ and
$\beta\colon[n]\twoheadrightarrow[q]$, and $y,z$ are nondegenerate.
For any order-preserving section $a\colon[p]\to[n]$ of $\alpha$,
\[
 y=a^*x=(\beta a)^*z.
\]
By Exercise~\ref{ex:3}, factor $\beta a$ as a surjection followed by
an injection. Nondegeneracy of $y$ forces the surjection to be an
identity, so $\beta a$ is injective and $p\le q$.
Interchanging the two expressions gives $q\le p$. Thus $p=q$, and
the increasing injection $\beta a\colon[p]\to[p]$ is the identity
for every section $a$ of $\alpha$.

Every $t\in[n]$ belongs to the image of some such section: choose
$t$ in its fiber and one element in each other fiber of $\alpha$;
these choices are automatically increasing. Hence
$\beta(t)=\alpha(t)$ for every $t$, so $\alpha=\beta$.
Applying any section now gives $y=z$.
\end{proof}
\end{exercise}
